\section*{Chapter \ref{ch:b1:consumer-choice} --- Consumer Choice}
% Numbers recomputed in code/bachelor-1/05-consumer-choice/exercises/solutions.py (tested).

\begin{solution}{exo:b1:consumer-choice:cd}
By \cref{prop:b1:consumer-choice:cd-demand}, $x_1 = 0.25 \times 100/2 = 12.5$
and $x_2 = 0.75 \times 100/5 = 15$. She spends $25$ on good 1 and $75$ on
good 2: shares of $25\%$ and $75\%$, the exponents.
\end{solution}

\begin{solution}{exo:b1:consumer-choice:read}
$3600/37935 = 9.5\%$ and $10308/159512 = 6.5\%$. The top fifth spends
$10308/3600 = 2.86$ times as many dollars on food at home. Spending rises
across the fifths, so food at home is \omterm{def:b1:consumer-choice:normal-inferior}{normal} there (taking prices to be the
same for both fifths); its share falls because spending on it rises less
than in proportion to the budget.
\end{solution}

\begin{solution}{exo:b1:consumer-choice:pivot}
Intercepts $100/2 = 50$ and $100/5 = 20$, slope $-2/5 = -0.4$. With $p_1 = 4$:
intercepts $25$ and $20$, slope $-0.8$. A bundle of the old line costs
$4x_1 + 5x_2 = (2x_1 + 5x_2) + 2x_1 = 100 + 2x_1$ at the new prices, more than
100 unless $x_1 = 0$: only $(0, 20)$ stays affordable.
\end{solution}

\begin{solution}{exo:b1:consumer-choice:normal}
Bread rises with income, so it is \omterm{def:b1:consumer-choice:normal-inferior}{normal} over the range; bus trips fall, so
they are inferior. With $\ln(40000/15000) = 0.981$: bread
$\ln(140/120)/0.981 = 0.154/0.981 = 0.157$; bus trips
$\ln(220/300)/0.981 = -0.310/0.981 = -0.316$.
\end{solution}

\begin{solution}{exo:b1:consumer-choice:min}
The optimum is at the corner of an L, $x_1 = 3x_2$. The budget
$p_1(3x_2) + p_2x_2 = m$ gives $x_2 = m/(3p_1 + p_2)$ and $x_1 = 3m/(3p_1 + p_2)$.
At $p_1 = 2$, $p_2 = 1$, $m = 70$: $x_2 = 70/7 = 10$ and $x_1 = 30$.
\end{solution}

\begin{solution}{exo:b1:consumer-choice:quasi}
The tangency $2/\sqrt{x_1} = p_1$ gives $x_1 = 4/p_1^2$, which costs
$p_1x_1 = 4/p_1$; then $x_2 = m - 4/p_1$. This is the optimum when
$m \ge 4/p_1$. Below that income she spends everything on good 1:
$x_1 = m/p_1$, $x_2 = 0$ (the Kuhn--Tucker condition $u_2 = 1 \le \lambda$ holds
because $\lambda = u_1/p_1 > 1$ there). At $p_1 = 0.5$: $x_1 = 16$, which costs 8.
With $m = 20$, $(x_1, x_2) = (16, 12)$; with $m = 5 < 8$, $(10, 0)$.
\end{solution}

\begin{solution}{exo:b1:consumer-choice:roy}
$\partial v/\partial p_1 = -a\,v/p_1$ and $\partial v/\partial m = v/m$, so
$-(\partial v/\partial p_1)/(\partial v/\partial m) = (av/p_1)(m/v) = am/p_1$,
the demand of \cref{prop:b1:consumer-choice:cd-demand}.
\end{solution}

\begin{solution}{exo:b1:consumer-choice:predict}
A $10\%$ rise predicts $0.31 \times 10\% = 3.1\%$ more food spending:
$4645 \times 0.031 = 144$ dalens a year. The prediction treats the
cross-section slope as the response of one household over time. That
fails if households with high consumption differ from the others in ways
that also move food spending (size, age, tastes): then the slope mixes the
effect of consumption with those differences.
\end{solution}

\begin{solution}{exo:b1:consumer-choice:warp-table}
Costs at each observation's prices:
at A's prices, $B$ costs 8 and $C$ costs 9, both $\le 10$, so A is directly
revealed preferred to B and to C. At B's prices, A costs 15 $> 14$ and C
costs $4 + 7 = 11 \le 14$, so B is directly revealed preferred to C. At C's
prices, A costs $5 + 10 = 15 \le 16$ and B costs $6 + 4 = 10 \le 16$, so C is
directly revealed preferred to A and to B. Pair A--B: \omterm{def:b1:consumer-choice:warp}{WARP} holds. Pair
A--C: each is revealed preferred to the other, a violation. Pair B--C: a
violation too.
\end{solution}

\begin{solution}{exo:b1:consumer-choice:prove}
Homogeneity: $B(tp, tm) = B(p, m)$ for $t > 0$, and a unique best bundle in
one set is the best bundle in the same set. Exhaustion: the proof of
\cref{prop:b1:consumer-choice:homogeneity} uses monotonicity only to show
that a bundle with a little more good 1 is better. Without it: take
$u = -(x_1 - 1)^2 - (x_2 - 1)^2$, prices 1 and income 10. The bundle $(1, 1)$
is the best of all bundles and costs 2, so the consumer spends 2 of her 10.
\end{solution}

\begin{solution}{exo:b1:consumer-choice:aggregate}
$X_1 = (a_1m_1 + a_2m_2)/p_1$. If $a_1 = a_2 = a$, $X_1 = a(m_1 + m_2)/p_1$,
which depends only on the total. Moving 10 from consumer 2 to consumer 1
changes $X_1$ by $(a_1 - a_2) \times 10/p_1 = 0.2 \times 10/2 = 1$ unit: the
consumer who spends a larger share on good 1 now has the money.
\end{solution}

\begin{solution}{exo:b1:consumer-choice:refute}
The hypothesis predicts the same budget shares in every household,
whatever its spending. A share that changes systematically with total
spending refutes it, and the opening chart shows such shares: food at home
from $9.5\%$ to $6.5\%$, shelter from $27.3\%$ to $19.4\%$ across the fifths of
2025. The hypothesis could be saved by letting the Cobb--Douglas share
differ between households, with poorer households happening to have higher
food shares. Then any pattern of shares across the fifths can be produced,
and the model predicts nothing about the chart; it can only be tested by
following households as their spending changes.
\end{solution}

\begin{solution}{pb:b1:consumer-choice:survey}
\textbf{1.} A: $12000/4 = 3000$ kg or $12000$ dalens of other goods; B: 5,000 kg
or 20,000; C: 10,000 kg or 40,000. Every line has slope $-4$: four dalens
of other goods per kilogram.

\textbf{2.} $12000/5 = 2400$ kg; slope $-5$.

\textbf{3.} At the new price it costs $5 \times 900 + 8400 = 12900 > 12000$: no.
An extra 900 dalens would make it affordable.

\textbf{4.} $4f + o = 4 \times 600 + 6000 = 8400$.

\textbf{5.} $5f + o = 5 \times 600 + 6000 = 9000$; the point $(600, 6000)$ satisfies
both equations. If D chose to sell food before the rise ($f < 600$), its old
bundle costs $5f + o = (4f + o) + f = 8400 + f < 9000$ at the new prices, so
it is still affordable with money to spare: D is better off. A net buyer of
food may be worse off.

\textbf{6.} Nothing: $1.1p_1x_1 + 1.1p_2x_2 \le 1.1m$ is the same set.

\textbf{7.} With prices of 1, \cref{prop:b1:consumer-choice:cd-demand} gives
$f = \beta c$: the share is $\beta$ in every household.

\textbf{8.} $\hat\beta = 0.199$. Poorest tenth: $0.199 \times 11639 = 2316$
against 3,726 in the survey. Richest: $0.199 \times 42736 = 8504$ against
5,594. The model is far too low at the bottom and far too high at the top.

\textbf{9.} $\ln f = \ln\beta + \ln c$: a slope of 1, against $0.31$ in the survey.

\textbf{10.} The same share, 0.199, in every tenth; the survey's share falls
from $0.320$ in the poorest tenth to $0.131$ in the richest.

\textbf{11.} Yes: food spending rises from 3,726 to 5,594 between the poorest
and the richest tenths. And it is a necessity: the elasticity $0.31$ is
below one.

\textbf{12.} Yes, if households with lower consumption happened to have
higher shares $\beta_h$. It would require tastes that vary with consumption
in just the way the survey shows. One cross-section cannot tell this apart
from common preferences that are not Cobb--Douglas: both give the same
cloud. Following the same households over time could.

\textbf{13.} Yes, in the same direction: the Cobb--Douglas model predicts a
flat share and an elasticity of 1. The American elasticity, $0.73$, is
closer to one than Brenland's.

\textbf{14.} Put $g = f - \gamma$; then $o = c - f = (c - \gamma) - g$, and the problem is
to maximise $\beta\ln g + (1 - \beta)\ln o$ subject to $g + o = c - \gamma$: a
Cobb--Douglas problem with income $c - \gamma$. So $g = \beta(c - \gamma)$, and
$f = \gamma + \beta(c - \gamma)$.

\textbf{15.} $f = (1 - \beta)\gamma + \beta c$: a straight line with intercept
$(1 - \beta)\gamma$ and slope $\beta$.

\textbf{16.} $\beta = 0.056$ and $(1 - \beta)\gamma = 3334$, so
$\gamma = 3334/0.944 = 3532$ dalens a year.

\textbf{17.} Poorest: $3334 + 0.056 \times 11639 = 3334 + 652 = 3986$, a share of
$3986/11639 = 0.342$ (survey 3,726 and 0.320). Richest:
$3334 + 0.056 \times 42736 = 3334 + 2393 = 5727$, a share of
$5727/42736 = 0.134$ (survey 5,594 and 0.131). Far closer than the
Cobb--Douglas predictions of question 8.

\textbf{18.} Poorest: $\eta = 652/3986 = 0.164$. Richest: $\eta = 2393/5727 = 0.418$.
The elasticity rises with consumption, while the log fit imposes $0.31$
everywhere.

\textbf{19.} $\eta = \beta c/\bigl((1 - \beta)\gamma + \beta c\bigr) < 1$ because
$(1 - \beta)\gamma > 0$. As $c$ grows, $\eta \to 1$: far above subsistence the
consumer behaves like a Cobb--Douglas one.

\textbf{20.} The model gives them no utility at all, since $\ln(f - \gamma)$ is not
defined for $f < \gamma$: one common $\gamma$ cannot describe 433 of the 2,000
households. The survey records household size, and larger households need
more food: letting $\gamma$ depend on size, $\gamma_h = \gamma_0 + \gamma_1 \times \text{size}_h$,
would let each household have its own subsistence quantity.

\textbf{21.} $1.1 \times 11639 = 12803$. Cobb--Douglas: food rises by $10\%$ too,
$1.1 \times 3726 = 4099$, and the share stays at $0.320$.

\textbf{22.} Food rises by $0.31 \times 10\% = 3.1\%$: $3726 \times 1.031 = 3842$, a share of
$3842/12803 = 0.300$.

\textbf{23.} From the line's own prediction: $3334 + 0.056 \times 12803 = 3334 + 717 = 4051$,
a share of $4051/12803 = 0.316$ (up 65 from 3,986). From the observed
spending: $3726 + 0.056 \times (12803 - 11639) = 3726 + 0.056 \times 1164 = 3726 + 65 = 3791$,
a share of $3791/12803 = 0.296$.

\textbf{24.} Rises of $4099 - 3726 = 373$ (Cobb--Douglas), $3842 - 3726 = 116$ (log
slope) and 65 (straight line). The Cobb--Douglas prediction is almost six
times the straight line's. Fitted to the same cross-section, the three
models disagree most about how the poorest households respond.

\textbf{25.} Observations of the same households before and after their
income changes (a panel), ideally when the change has a cause unrelated to
their tastes, such as a change in a transfer rule.

\textbf{26.} The subsistence food spending that fits the survey is
$\gamma = 3532$, about 3,530 dalens a year: the intercept of the straight-line
\omterm{def:b1:consumer-choice:engel-curve}{Engel curve}, divided by $1 - 0.056$.
\end{solution}
